% Generated draft: requires mathematical and source review.
\documentclass[11pt]{article}
\usepackage[margin=25mm]{geometry}
\usepackage{fontspec}
\setmainfont{Latin Modern Roman}
\usepackage{amsmath,amssymb,mathtools,unicode-math}
\setmathfont{Latin Modern Math}
\usepackage{enumitem,needspace,xurl,hyperref}
\setlength{\parindent}{0pt}
\setlength{\parskip}{5pt}
\setlength{\emergencystretch}{3em}
\newcommand{\N}{\mathbb{N}}
\newcommand{\Z}{\mathbb{Z}}
\newcommand{\Q}{\mathbb{Q}}
\newcommand{\R}{\mathbb{R}}
\newcommand{\C}{\mathbb{C}}
\newcommand{\F}{\mathbb{F}}
\newcommand{\A}{\mathbb{A}}
\newcommand{\PP}{\mathbb{P}}
\newcommand{\E}{\mathbb{E}}
\newcommand{\eps}{\varepsilon}
\newcommand{\dd}{\,\mathrm{d}}
\begin{document}
\section*{UnsolvedMath ID: 30005364}
\textbf{Source problem code:} catalogue entry
\subsection{Definitions and mathematical statement}
Let \(X\) be a smooth \(4\)-manifold containing an embedded nonseparating \(3\)-manifold \(Y\), and let \(\Sigma\) be the central surface of a trisection of \(X\). In the connected ambient case, nonseparation means that \(X\setminus Y\) is connected. The supplied statement does not explicitly specify compactness, boundary conditions, orientability, or connectedness hypotheses for these manifolds; these must be checked before fixing the precise trisection setting.

Write \(D_0,D_1,D_2\) for the three disk sets associated with the trisection. In the usual handlebody interpretation, a disk set consists of isotopy classes of essential simple closed curves on \(\Sigma\) that bound properly embedded disks in the corresponding three-dimensional handlebody. The supplied excerpt does not specify the ambient complex containing these sets, its edges, or the convention for measuring loop length. Its definition of trisection distance can therefore only be transcribed schematically as
\[
 d(\Sigma)=\min\bigl\{\ell(\gamma):\gamma\text{ is a loop in }D_0\cup D_1\cup D_2,\quad \gamma\cap D_i\ne\varnothing\text{ for }i=0,1,2\bigr\}.
\]
Here the meaning of loop, the length \(\ell\), and the structure on the union remain to be recovered from the cited report. This display does not select an ambient curve graph or another complex. The convention permitting distance zero also requires clarification.

\textbf{Problem 4.} Under the hypotheses and conventions intended in the cited report, use the topology of \(Y\) and the loop induced by \(Y\) in \(D_0\cup D_1\cup D_2\) to give a general upper bound for \(d(\Sigma)\). The supplied record does not define the induced-loop construction or identify which topological quantities of \(Y\) should enter the bound. Consequently, it does not determine a specific numerical inequality or a precise domain of admissible bounding data.

\textbf{Qualifications from the surrounding excerpt.} The source passage states that a nonseparating \(Y\) cannot avoid a surface isotopic to \(\Sigma\): it argues that avoidance would make the inclusion \(Y\hookrightarrow X\) factor through a \(2\)-complex, contradicting an asserted injection \(H_3(Y)\hookrightarrow H_3(X)\). It also describes removing intersections consisting only of trivial curves with an unspecified level surface \(h^{-1}(s)\), by ambient \(2\)-handle attachments producing a \(3\)-manifold \(Y'\) homologous to \(Y\) and disjoint from that surface. These are supplied contextual assertions, not established consequences of the incompletely specified hypotheses here; homology coefficients, the map \(h\), and the surgery conventions are absent.

The excerpt additionally asserts that the trisection is reducible if and only if \(d(\Sigma)=0\). Its separately numbered Conjecture 5 states that if a smooth \(4\)-manifold \(X\) contains a nonseparating \(3\)-sphere, then every trisection of \(X\) is reducible. This conjecture is contextual and is distinct from the requested upper-bound problem.
\subsection{Short English statement}
The problem asks for an upper bound on trisection distance using an embedded nonseparating \(3\)-manifold and its induced loop through the three disk sets. The supplied excerpt leaves the ambient complex, loop construction, length convention, and some manifold hypotheses unspecified, so this draft preserves the request while identifying what must be recovered from the report.
\subsection{Sources}
\begin{itemize}
\item Local UnsolvedMath registry, \texttt{data/unsolvedmath/problems.json}, permanent ID \(30005364\), title: \emph{Trisection Distance Bounds from Nonseparating Hypersurfaces}; exact source code \texttt{catalogue entry}. The supplied record identifies the target as Problem 4 and supplies no individual proposer.
\item \emph{Low-dimensional Topology} (2023), Oberwolfach Reports, pp.\ 151--213, DOI \texttt{10.4171/owr/2023/3}: \url{https://doi.org/10.4171/owr/2023/3}. This is the supplied primary source reference.
\item Koenig, \emph{3-manifolds lying in trisected 4-manifolds}: \url{https://arxiv.org/abs/1806.04870}. This is supplied as a literature-review reference; the record does not identify a particular result or passage supporting the requested bound.
\item Supplied literature assessment, checked exactly \texttt{2026-08-22}: the record labels the problem open and states that the report presents it after defining trisection distance, with no later primary result located supplying the requested general upper bound in the stated setting. This is supplied dated curation derived from the Oberwolfach record, not an independently verified literature assessment. The supplied background explicitly requests independent verification of time-sensitive status and cited claims.
\end{itemize}
\end{document}
